\capitulo{2}{Objetivos}

Para comenzar con esta memoria se definirá el conjunto de objetivos que se pretenden alcanzar con la realización de este Trabajo de Fin de Grado. 



\section*{Objetivos Generales}

\begin{itemize}
    \item El propósito principal es proporcionar una herramienta interactiva (aplicación web) que permita informar y concienciar sobre el impacto de los factores de riesgo en el desarrollo del melanoma maligno de piel, en el contexto geográfico español.
\end{itemize}

\section{Objetivos Técnicos}

\begin{enumerate}
    \item Diseñar y desarrollar una aplicación web interactiva utilizando el \textit{framework} \textit{Shiny}, capaz de integrar datos sanitarios y climáticos para visualizar información relevante sobre el \textit{melanoma} maligno de piel.

    \item Construir una herramienta de estimación orientativa del riesgo acumulado de padecer \textit{melanoma}, basada en variables personales y ambientales, usando técnicas heurísticas inspiradas en modelos de medicina preventiva.

    \item Integrar mapas interactivos geográficos mediante la librería \texttt{leaflet} que permitan visualizar los datos procesados y limpios de las respuestas a las APIs seleccionadas.

    \item Implementar un recomendador dinámico de medidas preventivas personalizadas en función del tipo de piel del usuario y las condiciones climáticas actuales, mostrando mensajes adaptados al usuario.
\end{enumerate}

\section{Objetivos Personales}

\begin{itemize}
    \item Aplicar los conocimientos adquiridos durante el grado en el procesamiento, limpieza y estructuración de datos, con el objetivo de transformarlos en información clara, visual y comprensible para el público general.

    \item Diseñar y desplegar una aplicación web funcional y accesible, orientada a la divulgación en salud pública, utilizando herramientas de código abierto como el \textit{framework} \textit{Shiny}.

    \item Consolidar y ampliar las competencias en programación con R, ciencia de datos y salud pública, integrando los conocimientos teóricos adquiridos en un proyecto práctico con utilidad real.

    \item Familiarizarse con nuevas tecnologías y enfoques informáticos, especialmente en el desarrollo de interfaces interactivas, aplicando principios del diseño web, experiencia de usuario y visualización de datos.

    \item Contribuir al ámbito de la epidemiología ambiental mediante el desarrollo de un recurso accesible y educativo que facilite la prevención y concienciación sobre el melanoma maligno de piel.
\end{itemize}







